\section{Formalization status and a concrete completion plan}
\label{sec:formalization}

The report contains complete written proofs of its stated restricted
mathematical results, with an explicit arithmetic induction audit. The
archive adds executable definitions, exact finite checks and a checked
part of the embedded proof infrastructure. Those categories have different
meanings; the following inventory makes each claim reviewable.

\subsection{What was checked}

The standalone Lean files were compiled with Lean 4.32.0, using only the
standard distribution. The compiler commit was
\begin{center}\small\code{8c9756b28d64dab099da31a4c09229a9e6a2ef35}.\end{center} The arithmetic kernel is a source-preserving extraction from
the pinned repository, with imports adapted as described in
Section~\ref{sec:object-iteration}. The build includes:
\begin{longtable}{@{}>{\raggedright\arraybackslash}p{.27\linewidth}>{\raggedright\arraybackslash}p{.35\linewidth}>{\raggedright\arraybackslash}p{.32\linewidth}@{}}
\toprule
File & Content & Verification meaning \\
\midrule\endhead
\code{ObjectKernel.lean} & Unchanged syntax, natural deduction and iPRA
declarations from three upstream modules & Kernel checked after import
adaptation; not new mathematical results \\
\code{ObjectIteration.lean} & Substitution lemmas and the general
primitive-recursion invariant rule & Actual embedded \code{Prf PRAAx}
proof transformation \\
\code{ObjectExample.lean} & Iterated successor equals its iteration count &
Closed embedded arithmetic derivation using the new rule \\
\code{NegativeSyntax.lean} & Negative translation, syntax laws, preservation
of QF, translation of all source axioms & Proved syntax facts and actual
embedded derivations of translated axioms; full proof translation pending \\
\code{ODETargets.lean} & Polynomial tables, Euler algorithm, solution name,
coefficient certificate and exact goal propositions & Definitions and target
types checked; the goal propositions are unproved \\
\code{SplitTargets.lean} & Integer polynomial algorithms, certificate,
bounds, enumeration and goal propositions & Definitions and target types
checked; the goal propositions are unproved \\
\code{Audit.lean} & Axiom reports and explicit target-type checks & Records
dependencies; checking a target's type does not establish it \\
\bottomrule
\end{longtable}

The audited new theorems use Lean's standard \code{propext} and
\code{Quot.sound} in the metatheory; some require only the former. The
reports contain no \code{sorryAx} or \code{Classical.choice} dependency.
Those metatheoretic principles do not add excluded middle, arbitrary
induction or choice to the embedded \code{PRAAx} proof calculus. The
source theory \code{ClassicalPRAAx} is an explicitly defined datatype of
axioms used as the translation input, not a new trusted Lean axiom.

The full repository native adapter was inspected against the current API
but was not built against the whole repository dependency closure. It is
clearly separated from the standalone build. No full object-theory proof
of the ODE construction, separated-branch bound or logical proof translator
is claimed.

\subsection{Definitions supplied for the analytic development}

\code{ODETargets.lean} uses finite coordinate functions
\code{Fin d -> Rat} for rational vectors. A monomial consists of a rational
coefficient and a finite exponent vector; a polynomial is a list of
monomials. Repeated exponent vectors are allowed. The polynomial-zero
checker adds the coefficients belonging to each exponent vector occurring
in the list and checks that all those sums vanish. Thus it handles
cancellation, and does not use equality of functions on an infinite
domain as a computational test.

The file defines the exact Euler-node recursion, rational polygon
interpolation, coefficient magnitude and Lipschitz bounds, the automatic
local horizon, the integer growth and Cauchy constants, and the reindexed
name. The predicate \code{Controlled} is precisely the finite-precision
solution contract in Definition~\ref{def:controlled-flow}.
\code{SameName} states agreement at arbitrarily accurate common
precisions. Its use is metatheoretic: each fixed evaluator becomes a PR
function symbol in the embedded version. The main target names are
\code{RefinementGoal}, \code{LocalExistenceGoal}, \code{UniquenessGoal}
and \code{ConservationGoal}.

The targets quantify over finite polynomial input data, and the uniform
algorithm evaluates that one finite grammar. This accommodates variable
dimension and degree without an unclocked interpreter for arbitrary PR
syntax.

\subsection{Dependency order for the main proofs}

Here is a path with independently reviewable endpoints, rather than a list
of unproved high-level axioms.

\paragraph{Stage 1: numerical representations.}
Implement natural encodings for integers, positive-denominator rationals,
finite arrays and polynomial tables. Prove the local ring and order laws,
finite-product absolute-value rules and positive-denominator cancellation
inside \code{Prf PRAAx}. Implement Boolean folds for bounded assertions.
The invariants in Appendix~\ref{sec:coding} describe the numerical
induction formulas. The existing repository may supply parts of this
library, but their actual proof endpoints must be imported explicitly.

\paragraph{Stage 2: reusable finite estimates.}
For ODEs, formalize coefficient bounds, rational Bernoulli, the finite
binomial estimate $(1+1/N)^N\le3$, and the variable-error Gronwall fold.
Each is an algebraic inequality with bounded index conditions. For split
branches, formalize \code{intFactorAbsLeGoal},
\code{polyEvalAbsLeGoal}, \code{polyDominanceGoal} and
\code{nonzeroPolyNoLargeRootGoal}. These are concrete named propositions
in \code{SplitTargets.lean}; their mathematical proofs are in
Section~\ref{sec:split-branches}.

The general rule of Theorem~\ref{thm:object-rule} discharges the outer
primitive-recursion induction after these local step proofs are supplied.
It is especially useful for folding bounded maxima, products, Euler tables
and accumulated errors. It does not replace the arithmetic in the local
step proofs.

\paragraph{Stage 3: the ODE invariant and refinement.}
Formalize the node bound
$\norm{x_{n,j}-x_0}\le jhM$, with the guard $j\le2^n$ encoded as a
Boolean or a quantifier-free implication. The paired-mesh state stores
both current nodes, parity of the fine-node index, and accumulated
error. Prove the recurrence in Theorem~\ref{thm:ode-refine}, apply the
finite Gronwall bound, and then extend from endpoints by affine
interpolation. A second finite fold telescopes refinements.
This proves \code{RefinementGoal} and the regularity component of
\code{LocalExistenceGoal}.

\paragraph{Stage 4: ODE meaning and conservation.}
The local-defect proof uses a finite partition of a rational interval by
mesh points, exact midpoint identities, and a rational triangle inequality.
It completes the controlled-solution contract. For uniqueness, fix the
requested accuracy first, compute $N,p$, and compare the finite table of
two candidates at that precision. For conservation, prove coefficient
checker soundness and the multivariate binomial remainder, then use one
accumulated-error invariant. Finally introduce the accuracy challenge
universally. No induction formula asserts equality of arbitrary real
functions.

\paragraph{Stage 5: complete integral-point lists.}
Prove \code{splitXBoundGoal} by coefficient dominance, and
\code{splitYBoundGoal} by the two cases $R(x)=0$ and $R(x)\ne0$.
Signed-interval membership and a rectangle scan give
\code{enumerateSolutionsIffGoal}. The numerical witnesses for a point in
the rectangle are $x+H$ and $y+K$; their bounds are already known.
The general family classification then needs only the integer unit
factors, $x\in\{-1,0,1\}$, and the parity fold described in its proof.

\paragraph{Stage 6: logical translation and native transport.}
The syntax and axiom layers of negative translation are already proved.
The next new lemma is the arithmetic equality-stability adapter, followed
by formula stability and induction on a complete source derivation.
The proof cases are fully spelled out in Section~\ref{sec:logic}.
After producing the desired \code{Prf PRAAx} theorem, apply the existing
\code{pra\_iff\_native} bridge. Its use is the final step, not an
alternative to Stages 1--5.

\subsection{Finite reference computations}

The exact-rational ODE program checks 5,342 assertions over exponential,
oscillator, nonlinear local-flow, Hermite-accumulator and circle-area
examples. It checks the finite inequalities and four coefficientwise
conservation certificates, and rejects an invalid certificate. The
nonlinear example uses small meshes because exact rational numerators and
denominators grow rapidly; the proof itself covers arbitrary precision.
No floating-point estimate is used as a proof oracle.

The integer program checks the complete certificate rectangle for the
mixed cubic, several further accepted certificates, exceptional fibers
where the remainder vanishes, coefficient expansion consistency, the
explicit family at $m=3,\ldots,12$ on stated finite boxes, and malformed
certificate rejection. The uniform classification and the absence of
solutions outside the certificate rectangle follow from the written
theorems, not from finite testing.

The source files and evidence record how to reproduce these runs. These
programs are reference implementations; neither is an iPRA proof checker.

\subsection{The difficulties that remain}

The framework makes the missing work more specific. It does not infer a
general height bound from finiteness or a cardinality bound, manufacture
Mordell--Weil quotient representatives, or establish the missing uniform
analytic estimates in prime-distribution packets. It does not decide
whether an arbitrary represented real is zero, choose singular Puiseux
branches from arbitrary Cauchy inputs, or infer a Cauchy modulus from finite
window stability. The source reports' warnings about these operations
remain applicable.

What is resolved here is narrower and concrete: an actual embedded
iteration rule, a written logical transfer mechanism with a checked axiom
layer, a constructed polynomial-flow interface for a stated analytic
unification problem, and a further explicit class of integral-point
problems with complete primitive-recursive solution lists.
